%; whizzy paragraph -pdf xpdf -latex ./whizzypdfptex.sh
%; whizzy-paragraph "^\\\\begin{frame}\\|\\\\emtext"
% latex beamer presentation.
% platex, latex-beamer でコンパイルすることを想定。 

%     Tokyo Debian Meeting resources
%     Copyright (C) 2012 Junichi Uekawa

%     This program is free software; you can redistribute it and/or modify
%     it under the terms of the GNU General Public License as published by
%     the Free Software Foundation; either version 2 of the License, or
%     (at your option) any later version.

%     This program is distributed in the hope that it will be useful,
%     but WITHOUT ANY WARRANTY; without even the implied warreanty of
%     MERCHANTABILITY or FITNESS FOR A PARTICULAR PURPOSE.  See the
%     GNU General Public License for more details.
%     You should have received a copy of the GNU General Public License
%     along with this program; if not, write to the Free Software
%     Foundation, Inc., 51 Franklin St, Fifth Floor, Boston, MA  02110-1301 USA

\documentclass[cjk,dvipdfmx,12pt]{beamer}
\usetheme{Tokyo}
\usepackage{monthlypresentation}

%  preview (shell-command (concat "evince " (replace-regexp-in-string "tex$" "pdf"(buffer-file-name)) "&")) 
%  presentation (shell-command (concat "xpdf -fullscreen " (replace-regexp-in-string "tex$" "pdf"(buffer-file-name)) "&"))
%  presentation (shell-command (concat "evince " (replace-regexp-in-string "tex$" "pdf"(buffer-file-name)) "&"))

%http://www.naney.org/diki/dk/hyperref.html
%日本語EUC系環境の時
\AtBeginDvi{\special{pdf:tounicode EUC-UCS2}}
%シフトJIS系環境の時
%\AtBeginDvi{\special{pdf:tounicode 90ms-RKSJ-UCS2}}

\newenvironment{commandlinesmall}%
{\VerbatimEnvironment
  \begin{Sbox}\begin{minipage}{1.0\hsize}\begin{fontsize}{8}{8} \begin{BVerbatim}}%
{\end{BVerbatim}\end{fontsize}\end{minipage}\end{Sbox}
  \setlength{\fboxsep}{8pt}
% start on a new paragraph

\vspace{6pt}% skip before
\fcolorbox{dancerdarkblue}{dancerlightblue}{\TheSbox}

\vspace{6pt}% skip after
}
%end of commandlinesmall

\title{東京エリアDebian勉強会}
\subtitle{第153回 2017年8月度}
\author{杉本 典充  dictoss@live.jp}
\date{2017年8月19日}
\logo{\includegraphics[width=8cm]{image200607/openlogo-light.eps}}

\begin{document}

\begin{frame}
\titlepage{}
\end{frame}

\begin{frame}{Agenda}
 \begin{minipage}[t]{0.45\hsize}
  \begin{itemize}
  \item 最近あったDebian関連のイベント報告
	\begin{itemize}
	\item 第152回東京エリアDebian勉強会
    \item Debian / Ubuntu ユーザーミートアップ in 札幌 2017.07
    \item OSC 2017 Hokkaido
	\end{itemize}
  \item 事前課題発表
  \end{itemize}
 \end{minipage}
 \begin{minipage}[t]{0.45\hsize}
  \begin{itemize}
   \item Debian Trivia Quiz
  \end{itemize}
 \end{minipage}
\end{frame}

\section{イベント報告}
\emtext{イベント報告}

\begin{frame}{第152回東京エリアDebian勉強会}
\begin{itemize}
\item 2017年6月18日に、サイボウズさんをお借りして開催しました
\item 参加者は20名でした
\item 「Debian Continuous Integration」(Charles Plessyさん）
\item 「Debian Policyドキュメント改良の試み」（発表者：やまねひでき）
  \item Debian 9 Stretch リリースパーティ
\end{itemize} 
\end{frame}

\begin{frame}{Debian / Ubuntu ユーザーミートアップ in 札幌 2017.07}
\begin{itemize}
\item 2017年7月14日に、インフィニットループさんをお借りして開催しました
\item 参加者は東京から2名、北海道から5名でした
\item 「lxcについて」（発表者：杉本典充さん）
\item 「Debian 9 Stretchのネットワークインターフェース名について」（発表者：吉野与志仁さん）
\item 北海道や札幌のOSS事情の情報を交換しながら、懇親を深めました
\end{itemize} 
\end{frame}

\begin{frame}{OSC 2017 Hokkaido}
\begin{itemize}
\item 2017年7月15日に、札幌コンベンションセンターさんにて開催しました
\item セミナー「Debian Updates」。主にDebian 9 Stretchの情報を説明しました。
\item ブース展示を行い、イベント来場者へdebianの広報活動や情報交換をしました。
\end{itemize} 
\end{frame}

\section{事前課題}
\emtext{事前課題}
{\footnotesize
 \begin{prework}{ dictoss }
  \begin{enumerate}
  \item $B;H$C$F$$$^$9(B
  \item wnpp$B$rD/$a$F$_$k(B
  \end{enumerate}
\end{prework}

\begin{prework}{ Haruki TSURUMOTO(tSU\_RooT) }
  \begin{enumerate}
  \item $B$^$@;H$C$?$3$H$,$"$j$^$;$s(B
  \item $B4IM}$7$F$$$k%Q%C%1!<%8$N%a%s%F%J%s%9$d2a5n$N%P%0Js9p$N8+JV$7$J$I(B
  \end{enumerate}
\end{prework}

\begin{prework}{ kenhys }
  \begin{enumerate}
  \item $B;H$C$F$$$^$9(B
  \item debexpo(mentors.d.n)$B$N(Bmaster$B$H(Blive$B%V%i%s%A$,$I$l$/$i$$0c$&$+3NG'$7$F$*$/(B
  \end{enumerate}
\end{prework}

\begin{prework}{ henrich }
  \begin{enumerate}
  \item $B;H$C$F$$$^$9(B
  \item $B869F$+(BOpenQA$B$,$i$_(B
  \end{enumerate}
\end{prework}

\begin{prework}{ takaswie }
  \begin{enumerate}
  \item $B$^$@;H$C$?$3$H$,$"$j$^$;$s(B
  \item aplay $B$N(Brewrite.
  \end{enumerate}
\end{prework}

\begin{prework}{ yy\_y\_ja\_jp }
  \begin{enumerate}
  \item $B;H$C$F$$$^$9(B
  \item DDTSS?
  \end{enumerate}
\end{prework}

\begin{prework}{ OsamuAoki }
  \begin{enumerate}
  \item $B;H$C$F$$$^$9(B
  \item uscan hack (DEBCONF17$B$NOC!K(B
  \end{enumerate}
\end{prework}

\begin{prework}{ koedoyoshida }
  \begin{enumerate}
  \item $B;H$C$F$$$^$9(B
  \item $BL$Dj(B
  \end{enumerate}
\end{prework}

\begin{prework}{ ysaito }
  \begin{enumerate}
  \item $B;H$C$F$$$^$9(B
  \item $B$&!<$s!#(Bgolang?
  \end{enumerate}
\end{prework}

}

\subsection{問題}
%; whizzy-master ../debianmeetingresume201311.tex
% $B0J>e$N@_Dj$r$7$F$$$k$?$a!"$3$N%U%!%$%k$G(B M-x whizzytex $B$9$k$H!"(Bwhizzytex$B$,MxMQ$G$-$^$9!#(B
%

\santaku
{quiz}
{answer A}
{answer B}
{answer C}
{B}
{desc}


\section{セミナー1}
\emtext{debconf17参加報告}

\section{Hack time}
\emtext{Hack time}

%\begin{frame}{成果を記入下さい!}
%\end{frame}
  
\section{今後のイベント}
\emtext{今後のイベント}
\begin{frame}{今後のイベント}
  \begin{itemize}
  \item 2017/9/16(土) 東京エリアDebian勉強会（予定）
  \item 2017/10/21(土) 東京エリアDebian勉強会（予定）
    \begin{itemize}
    \item 日時、会場、セミナー内容は調整中
    \end{itemize}
  \item 勉強会では発表者を随時募集中です。ぜひともよろしくお願いします。
  \end{itemize}
\end{frame}

\end{document}

;;; Local Variables: ***
;;; outline-regexp: "\\([ 	]*\\\\\\(documentstyle\\|documentclass\\|emtext\\|section\\|begin{frame}\\)\\*?[ 	]*[[{]\\|[]+\\)" ***
;;; End: ***
